% SPDX-License-Identifier: CC-BY-SA-4.0
% Session: Form parallelism to concurrency
% Exercise: Amdahl's law

\begingroup

\begin{exercice}[La loi d'Amdahl]
  \label{exo:introduction/amdahl}

  Une proportion $p=70\,\%$ d'un programme séquentiel peut être parallélisée. On dispose d'une machine avec $n=10$~c\oe urs.

  \begin{question} 
  \item Quelle est l'accélération (speedup) attendue ?
  \end{question}
  \begin{correction}
    $$ S(10)=\dfrac{1}{(1-p)+\dfrac{p}{10}} = \dfrac{1}{0,3+\dfrac{0,7}{10}} \approx 2,70. $$
  \end{correction}

  \begin{question} 
  \item On trouve dans le commerce une machine avec $n=20$~c\oe urs. Quelle est l'accélération attendue ?
  \end{question}
  \begin{correction}
    $$ S(20)=\dfrac{1}{(1-p)+\dfrac{p}{20}} = \dfrac{1}{0,3+\dfrac{0,7}{20}} \approx 2,99. $$
  \end{correction}

  \begin{question} 
  \item Quelle est l'accélération maximale théorique lorsque $n \to \infty$ ?
  \end{question}
  \begin{correction}
    $$ \lim_{n\to\infty} S(n) = \dfrac{1}{1-p} = \dfrac{1}{0,3} \approx 3,33. $$
  \end{correction}

  Un ingénieur pense qu'il existe un moyen d'optimiser la partie séquentielle pour la rendre deux fois plus rapide.

  \begin{question} 
  \item Quelle proportion du programme est parallélisable dans la version otimisée ? 
  \end{question}
  \begin{correction}
    On suppose maintenant que la partie séquentielle est deux fois plus rapide,
    mais que la partie non parallélisable est inchangée.
    On obtient donc : 
    $$T_1^{\mathrm{opt}} = \dfrac{T_1 \times (1-p)}{2} + T_1 \times p. $$

    La proportion parallélisable est donc :
    $$ p_{\mathrm{opt}} = \dfrac{T_1 \times p}{T_1^{\mathrm{opt}}} = \dfrac{2p}{1+p} \approx 82\%. $$
  \end{correction}
  
  \begin{question} 
  \item Quelle accélération obtient-on alors, dans la parallélisation avec $n = 10$~c\oe urs ?
  \end{question}
  \begin{correction}
    $$ S_{\mathrm{opt}}(10)=\dfrac{1}{(1-p_{\mathrm{opt}})+\dfrac{p_{\mathrm{opt}}}{10}} \approx 3,86. $$
  \end{correction}

  \begin{question} 
  \item À coût égal, vaut-il mieux choisir d'acheter la machine avec $n=20$~c\oe urs ou payer l'ingénieur pour optimiser le code ? 
  \end{question}
  \begin{correction}
    Il faut comparer l'accélération par rapport à $T_{1}$,
    de $T_{20}$, le temps d'exécution du programme initial sur la machine à $20$~c\oe urs,
    et de $T_{10}^{\mathrm{opt}}$, le temps d'exécution du programme optimisé sur la machine à $10$~c\oe urs.
    \begin{itemize}
    \item $\dfrac{T_{20}}{T_1} = S(20) \approx 2,99.$
    \item $\dfrac{T^{\mathrm{opt}}_{10}}{T_1} = \dfrac{1-p}{2} + \dfrac{p}{10} \approx 4,55.$
    \end{itemize}    

    L'optimisation de la partie séquentielle est nettement plus rentable.
  \end{correction}
  
\end{exercice}

\endgroup
\endinput
